% Integrated FF.43; review copy only, do not input separately.
\begin{claim}{FF.43}{Proved}\label{clm:FFmetabelianpointdetection}
Let \(E/B_p\) be finite and let
\(a,b\in E\). There is a proper
finite extension \(L/E\) and
\(\beta\in L\setminus E\) such
that equality of the two pairs
\[
 \begin{split}
 &\bigl(\kappa_L(a-\beta),
                  \kappa_L(1-a+\beta)\bigr),\\
 &\bigl(\kappa_L(b-\beta),
                  \kappa_L(1-b+\beta)\bigr)
 \end{split}
                                                        \tag{FF.43.1}
\]
implies \(a=b\). In fact every
proper finite extension and every
\(\beta\in L\setminus E\) work.

Consequently the following criterion
detects the actual point function of
a pinned map. Let \(X/B_p\) be
of finite type and let
\(f,g:X_{\acute et}\to\mathbb A^1_{B_p,\acute et}\)
be pinned morphisms over \(BG_{B_p}\).
For every finite \(L/B_p\) and
\(\beta\in L\), set
\[
 Y_\beta=\mathbb A^1_L\setminus\{\beta,\beta+1\},
 \qquad
 U_\beta=f_L^{-1}(Y_\beta)\cap g_L^{-1}(Y_\beta)
                         \subset X_L.
                                                        \tag{FF.43.2}
\]
Suppose the pullbacks by \(f_L\)
and \(g_L\) of the two Kummer
classes of \(T-\beta\) and
\(\beta+1-T\) agree after restriction
to \(U_\beta\). Then the two
actual geometric point functions
are equal, including their values
in \(K_p^{\rm ab}\).
This is a uniqueness criterion for
point functions; it neither constructs
the comparison morphism nor asserts
that equality of point functions
alone identifies arbitrary topos maps.
\end{claim}
\begin{proof}
Every finite extension of \(B_p\)
has a nontrivial absolute Galois
group, by the free-group argument
in Claim~\ref{clm:FFmetabeliantest}.
Thus a proper finite extension
\(L/E\) exists. If
\(\beta\in L\setminus E\), both
\(a-\beta\) and \(b-\beta\)
lie outside \(E\), hence outside
\(B_p\); neither equals zero or
one. The first value-detection
assertion of
Claim~\ref{clm:FFmetabeliankernel},
applied over \(L\), proves the
first conclusion.

For the criterion, take any
\(x\in X(\overline{\mathbb Q})\).
Choose a finite \(E/B_p\) over
which \(x\) is defined. Pinning
and the retained relative structure
give actual values
\(a=f(x),b=g(x)\in E\), by
\citep[Proposition~3.6]{FF_CHW2024}.
Choose \(L/E\) and \(\beta\)
as above. The point \(x_L\)
lies in both inverse-image opens
in (FF.43.2), since neither \(a\)
nor \(b\) is \(\beta\) or
\(\beta+1\). Evaluating the two
assumed cohomology equalities at
\(x_L\) gives (FF.43.1), so
\(f(x)=g(x)\). This argument
works at every geometric point
and makes no exception for values
in \(K_p^{\rm ab}\).
\end{proof}

The criterion uses genuine operations
available for pinned morphisms:
finite scalar extension, inverse images
of Zariski opens, and pullback and
evaluation of Kummer classes. Its
remaining hypothesis is substantial.
The visible nonsplit functional in
Claim~\ref{clm:FFmetabelianbarrier}
does not by itself produce compatible
equalities for these translated pairs,
or a scheme morphism with which to
compare them. Deriving those equalities
is the unresolved existence step.
